% The positive semidefinite cone in Sym_2(R), and its horizontal shadow.
%
% COORDINATES.  For A symmetric with entries A_00, A_01 = A_10, A_11 put
%   x_1 = (A_00 - A_11)/2,   x_2 = A_01,   x_3 = tau = tr(A)/2 .
% Then det A = x_3^2 - x_1^2 - x_2^2 and the two eigenvalues are
% x_3 +- sqrt(x_1^2 + x_2^2), so {A psd 0} is the CIRCULAR cone
% x_3 >= sqrt(x_1^2 + x_2^2) of half-angle 45 degrees about the x_3-axis, and
% its boundary is the rank-one locus.  A rank-one atom g g^T with
% g = r(cos th, sin th) has
%   x_1 = (r^2/2) cos 2th,   x_2 = (r^2/2) sin 2th,   x_3 = r^2/2 ,
% hence sits ON the boundary, at height half the leverage, at DOUBLE the
% argument; g and -g give the same point.  I_2 is the axis point (0,0,1).
% The chart is conformal: |A|_F^2 = 2(x_1^2 + x_2^2 + x_3^2), so coordinates to
% matrices is sqrt2 times a Euclidean isometry, the three coordinate directions
% I, diag(1,-1) and [[0,1],[1,0]] being Frobenius-orthogonal of norm sqrt2
% apiece.  The 45-degree half-angle is therefore the Frobenius one.
%
% PANEL (a): orthographic projection of R^3 with azimuth 75 and elevation 26
% degrees, scaled by 2.6cm:
%   X = 2.6(-0.965926 x_1 + 0.258819 x_2)
%   Y = 2.6(-0.113459 x_1 - 0.423434 x_2 + 0.898794 x_3)
% The circle of radius R at height tau projects to the ellipse centred at
% (0, 2.6*0.898794*tau) with semi-axes 2.6*R and 2.6*0.438371*R; at
% R = tau = 1 that is centre (0,2.3369), semi-axes 2.6 and 1.1398.  Its lower
% half is the near half and is drawn solid, its upper half hidden.  The
% silhouette rulings touch it where cos(phi - 75) = tan(26), i.e. at
% phi - 75 = +-60.8083 degrees, screen points (+-2.2698, 1.7810); they are
% carried on to tau = 1.14 so that the cone visibly continues past the slice.
%
% PANEL (b) is the plane of the atoms, face-on and unprojected, at 1.35cm per
% unit, with the first coordinate to the right.  g_c is dashed at length
% sqrt2, its square eta_c solid at length 2, and the dot at the midpoint of
% eta_c is the shadow of g_c g_c^T, at radius ell_c/2 = 1 on the dotted circle
% that is the shadow of the rim.
%
% THE DESIGN is the rank-two trine, at theta_c = 45, 165, 285 degrees:
%   g_0 = (1, 1),   g_1 = (-(1+sqrt3)/2, (sqrt3-1)/2),
%   g_2 = ((sqrt3-1)/2, -(1+sqrt3)/2),
% all of leverage 2, all of weight 1/3, with sum_c g_c g_c^T = 3 I_2.  Its
% squares are eta_0 = (0,2), eta_1 = (sqrt3,-1), eta_2 = (-sqrt3,-1), summing
% to 0, and its three atoms sit on the circle at height 1 at the plane angles
% phi_c = 2 theta_c = 90, 330, 210 degrees, that is at psi_c = phi_c - 75 =
% 15, 255, 135 from the view direction, so only the first is on the near half.
\begin{tikzpicture}[
  x=1cm, y=1cm,
  vec/.style={-{Latex[length=2mm,width=1.5mm]}, line width=0.75pt},
  gvec/.style={-{Latex[length=2mm,width=1.5mm]}, line width=0.75pt, black!55,
               dashed, dash pattern=on 1.6pt off 1.4pt},
  edge/.style={line width=0.45pt, black!65},
  hidden/.style={line width=0.45pt, black!45, dashed,
                 dash pattern=on 1.4pt off 1.2pt},
  frame/.style={densely dotted, line width=0.3pt, black!35},
  chord/.style={line width=0.45pt, black!58},
  note/.style={font=\scriptsize, align=center, inner sep=1.5pt},
  every node/.style={font=\small}]

% =====================================================================
% (a)  the cone, its rim at tau = 1, and the trine on that rim
% =====================================================================

% ---- the coordinate frame ------------------------------------------
\draw[frame] (0,-0.45) -- (0,3.95);
\node[note, anchor=west]  at ( 0.10, 3.95) {$\tau=\tfrac12\tr A$};
\draw[frame] (0,0) -- (-1.8836,-0.2212);      % x_1 = (A_00 - A_11)/2, 0.75
\node[note, anchor=north] at (-1.8836,-0.30) {$\tfrac12(A_{00}-A_{11})$};
\draw[frame] (0,0) -- ( 0.5047,-0.8257);      % x_2 = A_01, 0.75
\node[note, anchor=north] at ( 0.5047,-0.90) {$A_{01}$};

% ---- the cone: two silhouette rulings and the rim at tau = 1 --------
\draw[edge] (2.5876,2.0303) -- (0,0) -- (-2.5876,2.0303);
\draw[edge] ( 2.6,2.3369)
  arc[start angle=0,   end angle=-180, x radius=2.6, y radius=1.1398];
\draw[hidden] (-2.6,2.3369)
  arc[start angle=180, end angle=0,    x radius=2.6, y radius=1.1398];
\node[note, anchor=east] at (-1.30,0.85) {$\rk A=1$};

% ---- the rank-one rays carrying the two rear atoms, drawn as rulings
% on the far wall; the near wall is described by the silhouette instead.
\draw[frame] (0,0) -- (-2.5114,2.6319);
\draw[frame] (0,0) -- ( 1.8385,3.1428);

% ---- Parseval: the axis point is the centroid of the three ----------
\draw[chord] ( 0.6729,1.2359) -- (0,2.3369) -- (-2.5114,2.6319);
\draw[chord] (0,2.3369) -- (1.8385,3.1428);

\fill (0,0) circle (1.1pt);
\fill ( 0.6729,1.2359) circle (1.7pt);        % g_0, psi =  15, front
\fill (-2.5114,2.6319) circle (1.7pt);        % g_1, psi = 255, behind
\fill ( 1.8385,3.1428) circle (1.7pt);        % g_2, psi = 135, behind
\fill (0,2.3369) circle (1.7pt);              % I_2
\node[anchor=north]      at ( 0.55, 1.10) {$g_0^{\vphantom{\mathsf{T}}}g_0\tp$};
\node[anchor=south east] at (-2.56, 2.72) {$g_1^{\vphantom{\mathsf{T}}}g_1\tp$};
\node[anchor=south west] at ( 1.90, 3.21) {$g_2^{\vphantom{\mathsf{T}}}g_2\tp$};
\node[anchor=south west] at ( 0.12, 2.42) {$I_2$};

\node[note, anchor=north] at (0,-1.55)
  {$\sum_c t_c\,g_c^{\vphantom{\mathsf{T}}}g_c\tp=I_2$,\quad
   $t_c=\tfrac13$,\quad $\ell_c=2$\\[2pt]
   $\dim\mathrm{Sym}_2=3$,\quad $\win(2)=4$};

% =====================================================================
% (b)  the shadow: the atoms, their squares, and the doubled argument
% =====================================================================
\begin{scope}[xshift=6.95cm, yshift=0.85cm]

\draw[frame] (-2.95,0) -- (2.95,0);
\draw[frame] (0,0) circle (1.35);             % the shadow of the rim
\draw[frame] (1.35,-0.09) -- (1.35,0.09);
\draw[frame] ( 2.7,-0.09) -- ( 2.7,0.09);
\node[note, anchor=north] at (1.35,-0.11) {$1$};
\node[note, anchor=north] at ( 2.7,-0.11) {$2$};

% ---- theta_0 = 45 degrees, and then theta_0 again --------------------
\draw[chord] (1.48,0) arc[start angle=0, end angle=90, radius=1.48];
\node[note] at (1.6076,0.6658) {$\theta_0$};
\node[note] at (0.6658,1.6076) {$\theta_0$};

% ---- the three atoms, at length sqrt2 --------------------------------
\draw[gvec] (0,0) -- ( 1.3500, 1.3500);
\draw[gvec] (0,0) -- (-1.8441, 0.4941);
\draw[gvec] (0,0) -- ( 0.4941,-1.8441);
\node[note, anchor=south west] at ( 1.40, 1.40) {$g_0$};
\node[note, anchor=south east] at (-1.89, 0.54) {$g_1$};
\node[note, anchor=north west] at ( 0.55,-1.89) {$g_2$};

% ---- their squares, at length 2, the shadow at the midpoint ----------
\draw[vec] (0,0) -- ( 0.0000, 2.7000);
\draw[vec] (0,0) -- ( 2.3383,-1.3500);
\draw[vec] (0,0) -- (-2.3383,-1.3500);
\fill ( 0.0000, 1.3500) circle (1.7pt);
\fill ( 1.1691,-0.6750) circle (1.7pt);
\fill (-1.1691,-0.6750) circle (1.7pt);
\fill (0,0) circle (1.1pt);
\node[anchor=south] at ( 0.00, 2.76) {$\eta_0$};
\node[anchor=west]  at ( 2.41,-1.38) {$\eta_1$};
\node[anchor=east]  at (-2.41,-1.38) {$\eta_2$};

\node[note, anchor=north] at (0,-2.40)
  {$\lVert\eta_c\rVert=\ell_c=2$,\quad $\sum_c t_c\,\eta_c=0$};

\end{scope}
\end{tikzpicture}
